\documentclass[aspectratio=1610]{beamer}

\title{Potpuno homomorfna enkripcija}
\author{Tijana Jak\v si\' c}
\date{\today}

\usetheme[sectionpage=none]{wildcat}

\begin{document}

\begin{frame}
\titlepage
\end{frame}


\begin{frame}{Sadržaj}
\begingroup
\setbeamertemplate{sections/subsections in toc}[square]
\tableofcontents
\endgroup
\end{frame}

\section{Uvod}
% ---------------- MOTIVACIJA ----------------
\begin{frame}{\v Sta \v zelimo od homomorfne enkripcije?}
    \begin{itemize}
        \item Delegiranje vršenja operacija nad podacima, bez kompromitovanja privatnosti samih podataka
        \item Klijent delegira izračunavanje nad privatnim podacima serveru kome ne veruje
        \item Podatke šalje enkriptovane, server homomorfno vrši operacije nad enkriptovanim podacima, i šalje rezultat klijentu, koji klijent treba samo da dešifruje.
        \item Ideju za\v zeleli Rivest, Adleman i Dertouzos 1978, prvi predlog implementacije Gentri 2009.
    \end{itemize}
\end{frame}

\subsection{Multiplikativni homomorfizam}
% ---------------- MULTIPLIKATIVNI HOMOMORFIZAM RSA I ELGAMAL ----------------
\begin{frame}{Multiplikativni homomorfizam RSA}
    \begin{itemize}
        \item poruke $m_1$ i $m_2$ \v sifrovane sa $c_1 = m_1^e\pmod{N}$ i $c_2 = m_2^e\pmod{N}$
        \item $c_1\cdot c_2 = m_1^e\cdot m_2^e = (m_1\cdot m_2)^e \pmod{N}$
    \end{itemize}
\end{frame}

\begin{frame}{Multiplikativni homomorfizam ElGamal}
    \begin{itemize}
        \item poruke $m_1$ i $m_2$ \v sifrovane sa $c_1 = (g^{r_1},m_1\cdot h^{r_1})$ i $c_2 = (g^{r_2},m_2\cdot h^{r_2})$
        \item $c_1\odot c_2 = (g^{r_1}\cdot g^{r_2}, m_1\cdot h^{r_1}\cdot m_2\cdot h^{r_2}) = (g^{r_1 + r_2}, m_1\cdot m_2\cdot h^{r_1+r_2})$
    \end{itemize}
\end{frame}

\section{Notacija i definicije}
% ---------------- DEFINICIJA ----------------
\begin{frame}{Osnovni pojmovi}
    \begin{itemize}
        \item prsten $R(+, \cdot)$
        \item Aritmeti\v cka funkcija je kompozicija $+$ i $\cdot$
    \end{itemize}
    $$\mathrm{D}(\mathrm{E}(m, pk), sk) = m $$
    \begin{pblock}{Evaluacioni algoritam}
        Poruke $m_1, ..., m_n\in R$, njihovi \v sifrati $c_i = \mathrm{E}(m_i, pk)$, $i\in \{1, ..., n\}$ i aritmeti\v cka $f:R^n\rightarrow R$. \textit{Evaluacioni algoritam} je $\mathrm{Eval}$ takav da
        $$c_f = \mathrm{Eval}(f, (c_1, ..., c_n), ek)$$
        $$\mathrm{D}(c_f, sk) = f(m_1, ..., m_n)$$
    \end{pblock}
    \begin{itemize}
        \item $c = \mathrm{E}(f (m_1 , ..., m_n ), pk)$ i $c_f$ dešifruju u isti broj, ali su drugačiji u konstrukciji i u opštem slučaju $c \neq c_f$
    \end{itemize}
\end{frame}

\begin{frame}
    $$\mathrm{D}(\mathrm{E}(m_1, pk) \oplus \mathrm{E}(m_2, pk), sk) = m_1 + m_2$$
    $$\mathrm{D}(\mathrm{E}(m_1, pk) \odot \mathrm{E}(m_2, pk), sk) = m_1 \cdot m_2$$
    \begin{pblock}{SHE}
        Enkripciona \v sema je \textit{delimi\v cno homomorfna} ako postoji ograni\v cena klasa „dovoljno jednostavnih” aritmeti\v ckih funkcija $\mathcal{F}$ takva da za proizvoljnu $f\in \mathcal{F}$ proizvoljnog reda $n$, proizvoljne poruke $m_1, ..., m_n\in R$ i njihove \v sifrate $c_i = E(m_i, pk)$, $i\in \{1, ..., n\}$ va\v zi
    $$\mathrm{D}(\mathrm{Eval}(f, (c_1, ..., c_n), ek)) = f(m_1, ..., m_n)$$
    \end{pblock}
\end{frame}

\begin{frame}
    \begin{pblock}{FHE}
        Enkripciona \v sema je \textit{potpuno homomorfna} ako za proizvoljnu $f\in \mathcal{F}$ proizvoljnog reda $n$, proizvoljne poruke $m_1, ..., m_n\in R$ i njihove \v sifrate $c_i = E(m_i, pk)$, $i\in \{1, ..., n\}$ va\v zi
        $$\mathrm{D}(\mathrm{Eval}(f, (c_1, ..., c_n), ek)) = f(m_1, ..., m_n)$$
    \end{pblock}
\end{frame}

\subsection{Od SHE do FHE}
% ---------------- OD SHE DO FHE ----------------
\begin{frame}{Od SHE do FHE}
\begin{itemize}
    \item Bootstraping, proces smanjivanja \v suma
    \item Evaluacija u FHE se, u svim poznatim \v semama, vr\v si ponavljanjem slede\' ca dva koraka:
    \begin{enumerate}
        \item  Izvr\v siti dozvoljen broj operacija polaznom SHE \v semom
        \item  Pomo\' cu butstrepinga smanjiti \v sum 
    \end{enumerate}
\end{itemize}
\end{frame}

\section{Primer homomorfne enkripcije: DGHV}
% ---------------- UVODNI PRIMER ----------------
\begin{frame}{Uvodni primer}
\begin{itemize}
    \item Simetri\v cna enkripciona \v sema 
    \item \textbf{Generisanje ključa}: ključ je slučajan neparan broj $p$.
    \item \textbf{\v Sifrovanje}: $m \in \{0, 1\}$ se enkriptuje sa $$c=m + 2\cdot r + p\cdot q$$ gde je $r$ „mali” \v sum, a $q$ „veliki” slu\v cajan ceo broj
    \item \textbf{De\v sifrovanje}: $c$ se dekriptuje sa $$(c \pmod{p})\pmod{2}$$
\end{itemize}
\end{frame}

\begin{frame}{Aritmeti\v cke operacije}
\begin{itemize}
    \item Neka je $m_1$ \v sifrovano sa $c_1=m_1 + 2\cdot r_1 + p\cdot q_1$ i $m_2$ sa $c_2=m_2 + 2\cdot r_2 + p\cdot q_2$
    \item \textbf{Sabiranje}: $c_1 + c_2 = (m_1 + m_2) + 2\cdot (r_1 + r_2) + p\cdot (q_1 + q_2)$
    \item \textbf{Mno\v zenje}: $c_1\cdot c_2 = m_1\cdot m_2 + 2\cdot (r_2\cdot m_1 + r_1\cdot m_2 + 2\cdot r_1\cdot r_2) + p(q_1m_2 + 2q_1r_2 + 2q_1q_2 + 2q_2r_1 + q_2m_1)$
\end{itemize}
\end{frame}

% ---------------- DGHV ----------------
\begin{frame}{DGHV}
\begin{itemize}
    \item Javnim klju\v cem sakrivamo $p$
    \item \textbf{Generisanje ključa}: Tajni ključ je slučajan neparan $p$, javni je niz $(x_0, ..., x_n)$ td. $x_0$ najve\' ci i neparan, i za sve $i\in {1, ..., n}$ va\v zi $x_i = p\cdot q_i + r_i $ gde su 
$q_i, r_i$ slu\v cajni celi brojevi, s tim \v sto $r_0$ mora da bude paran.
    \item \textbf{\v Sifrovanje}: $m \in \{0, 1\}$ se \v sifruje u \[c = m + 2r + 2\sum_{i\in S}x_i \pmod{x_0}\] gde je $r$ slu\v cajan ceo broj, a  $S$ slu\v cajan podskup od $\{1, ..., n\}$.
    \item \textbf{De\v sifrovanje}: $c$ se de\v sifruje sa \[(c \pmod{p}) \pmod{2}\]
\end{itemize}
\end{frame}

\begin{frame}{De\v sifrovanje radi}
\begin{itemize}
    \item Od $c = m + 2r + 2\sum_{i\in S}x_i \pmod{x_0}$ se sre\dj ivanjem operacija do\dj e do oblika $c = m + 2R + pQ \pmod{x_0}$
    \item Da li smeta $\pmod{x_0}$?
    \item Postoji $k$ td. $m + 2R + pQ = c + kx_0$
    \item Uvo\dj enjem ovoga i sre\dj ivanjem operacija opet dolazimo do oblika $c = m + 2R' + pQ'$
\end{itemize}
\end{frame}

\begin{frame}{Aritmeti\v cke operacije}
\begin{itemize}
    \item Neka je $m_1$ \v sifrovano sa $c_1=m_1 + 2\cdot r_1 + p\cdot q_1$ i $m_2$ sa $c_2=m_2 + 2\cdot r_2 + p\cdot q_2$
    \item \textbf{Sabiranje}: $c_1 + c_2 = (m_1 + m_2) + 2\cdot (r_1 + r_2) + p\cdot (q_1 + q_2)$
    \item \textbf{Mno\v zenje}: $c_1\cdot c_2 = m_1\cdot m_2 + 2\cdot (r_2\cdot m_1 + r_1\cdot m_2 + 2\cdot r_1\cdot r_2) + p(q_1m_2 + 2q_1r_2 + 2q_1q_2 + 2q_2r_1 + q_2m_1)$
\end{itemize}
\end{frame}

\section{Bootstraping}
% ---------------- BOOTSTRAPING ----------------
\begin{frame}{Bootstraping}
    \begin{itemize}
        \item Od \v sifrata sa puno \v suma do \v sifrata iste poruke sa malo \v suma
        \item Algoritam de\v sifrovanja $\mathrm{D}(c, sk)$ uklanja \v sum
        \item I to je samo algoritam koji se mo\v ze sprovesti homomorfno u \v sifrovanom domenu
        \item Evaluacioni klju\v c je \v sifrat tajnog klju\v ca
        \item Pretostavka o kru\v znoj bezbednosti
    \end{itemize}
\end{frame}

\section{Primene homomorfne enkripcije}
% ---------------- PRIMENE ----------------
\begin{frame}{Primene}
    \begin{itemize}
        \item Cloud computing
        \item Bezbedna kolaboracija
        \item Ma\v sinsko u\v cenje
        \item Internet stvari
    \end{itemize}
\end{frame}

\section{Otvorena pitanja}
% ---------------- POLJA ZA NAPREDAK ----------------
\begin{frame}{Otvorena pitanja}
\begin{itemize}
    \item Kreirati efikasniji kriptosistem (\v sto manje overhead-a za bootstraping)
    \item FHE zasnovan na jo\v s nekoj primitivi osim re\v setki?
\end{itemize}
\end{frame}

\section{Zaklju\v cak}
\begin{frame}{Zaklju\v cak}
\begin{itemize}
    \item Homomorfna enkripcija omogu\' cava izra\v cunavanje u enkriptovanom domenu
    \item Poznate FHE \v seme se dobijaju od SHE \v sema procesom bootstrapinga
    \item Uveli smo \v semu DGHV kao SHE
    \item Klju\v cna ideja bootstrapinga je vr\v senje algoritma dekripcije homomorfno u \v sifrovanom domenu
    \item Naveli smo primene i otvorena pitanja
\end{itemize}
\end{frame}

\section{Pitanja}
\begin{frame}{Pitanja}
\begin{itemize}
    \item \v Sta je homomorfna enkripcija?
    \item \v Sta je bootstraping u homomorfnoj enkripciji?
    \item Navesti neku primenu homomorfne enkripcije.
\end{itemize}
\end{frame}

\standout{Hvala na pažnji!}

\end{document}
